\chapter{Matroids}

\begin{center}\itshape
Let $C$ be a finite set of elements, and let us suppose that certain subsets are called independent.\\
\upshape(Hassler Whitney, \emph{On the abstract properties of linear dependence}, 1935)
\end{center}

In 1926 the electricity company of western Moravia asked a young mathematician in Brno, Otakar Bor\r{u}vka, a practical
question: how should the towns of the region be connected to the network so that the total length of cable is as small
as possible? Bor\r{u}vka found the answer, and it was almost insultingly simple. Look at the possible connections from the
shortest to the longest, and lay each cable unless it would close a loop. Nothing cleverer is needed: the greedy choice,
never revised, is optimal. The same algorithm was rediscovered by Jarn\'\i k (1930), Kruskal (1956), Prim (1957) and Dijkstra (1959).

Greedy algorithms are usually wrong. A traveller who always goes to the nearest unvisited city can end up with a tour much longer
than the best one; a thief who fills a knapsack with the most valuable objects per kilogram may leave room for nothing
else. Why does greed work for networks of cables? In 1935 Hassler Whitney, at Harvard, noticed that the forests of a graph and
the linearly independent sets of vectors obey the same three rules, and he called any structure obeying them a
\emph{matroid}. (Takeo Nakasawa in Japan found the same idea the same year, in papers that were forgotten for half a century.)
Thirty-five years later Jack Edmonds proved the theorem that this chapter is built around: the greedy algorithm is optimal for
every choice of weights \emph{exactly} when the structure is a matroid.

Matroids are an abstraction of independence: of vectors, of edges, of matchable vertices, of points in general position.
They explain why the bases of a vector space all have the same size and why spanning trees of a graph all have the same number
of edges, by one argument. They have a duality that turns cycles into cuts and, for plane graphs, a graph into its dual map.
And some of them are strange: the Fano plane is a matroid that can be realised by vectors only over fields where $1+1=0$.

\begin{think}
Seven towns, and a table of the costs of connecting each pair. Kruskal's rule picks the cheapest connections one by one,
avoiding loops. Try to build a table of costs where this rule fails. Now do the same for the rule ``visit the nearest city next''
for a round trip through all towns.
\end{think}

\section{Independence}

\begin{definition}\label{def:matroid}
A \emph{matroid} $M=(E,\mathcal I)$ is a finite set $E$, the \emph{ground set}, with a family $\mathcal I$ of subsets, called \emph{independent}, such that
\begin{enumerate}[label=(I\arabic*)]
\item $\varnothing\in\mathcal I$;
\item if $B\in\mathcal I$ and $A\subseteq B$, then $A\in\mathcal I$;
\item (\emph{augmentation}) if $A,B\in\mathcal I$ and $|A|<|B|$, there is $x\in B\setminus A$ such that $A\cup\{x\}\in\mathcal I$.
\end{enumerate}
Subsets of $E$ that are not independent are \emph{dependent}.
\end{definition}

We write $A+x$ for $A\cup\{x\}$ and $A-x$ for $A\setminus\{x\}$. Axiom (I3) is the heart: a small independent set can always learn from a larger one.

\begin{proposition}[the basic examples]\label{prop:examples}
The following are matroids.
\begin{enumerate}[label=(\alph*)]
\item \emph{Uniform}: $E$ with $n$ elements and $\mathcal I$ the subsets of size at most $k$; written $U_{k,n}$.
\item \emph{Linear}: $E$ indexes a family of vectors $(v_e)_{e\in E}$ in a vector space over a field $F$, and $A$ is independent when the vectors $v_e$, $e\in A$, are
distinct and linearly independent. A matroid isomorphic to one of these is \emph{representable} over $F$.
\item \emph{Graphic}: $E$ is the edge set of a graph $G$, and $A$ is independent when it contains no cycle (it is a forest). Written $M(G)$, the \emph{cycle matroid} of $G$.
\item \emph{Partition}: $E=E_1\sqcup\dots\sqcup E_k$ with capacities $c_i\ge0$, and $A$ is independent when $|A\cap E_i|\le c_i$ for all $i$.
\item \emph{Transversal}: $S_1,\dots,S_m$ are subsets of $E$, and $A$ is independent when it is a \emph{partial transversal}: its elements can be assigned to distinct indices $i$ with
each element in its $S_i$.
\end{enumerate}
\end{proposition}
\begin{proof}
In each case (I1) and (I2) are clear, and (I3) is the point.

(a) and (d): add to $A$ an element of $B$ in a part where $A$ is below capacity; such a part exists, since $|B\cap E_i|>|A\cap E_i|$ for some $i$, and then $|A\cap E_i|<|B\cap E_i|\le c_i$.

(b) If $A+x$ were dependent for every $x\in B\setminus A$, then every $v_x$, $x\in B$, would lie in the span of $(v_a)_{a\in A}$ (for $x\in B\setminus A$ because $(v_a)_{a\in A}$ is independent and
$(v_a)_{a\in A+x}$ is not, or $v_x$ repeats a vector of $A$). So $|B|$ independent vectors would lie in a space of dimension $|A|<|B|$.

(c) A forest with edge set $A$ on the $n$ vertices of $G$ has $n-|A|$ components (each tree component has one edge fewer than vertices). If $|B|>|A|$, the forest $B$ has fewer components than
$A$, so some edge $x$ of $B$ joins two different components of $A$; otherwise each component of $B$ would lie inside a component of $A$, and $B$ would have at least as many components. Then $A+x$ is
still a forest, and $x\notin A$.

(e) Consider the bipartite graph with $E$ on one side, the indices $1,\dots,m$ on the other, and $e$ joined to $i$ when $e\in S_i$. A set is a partial transversal when some matching covers it. Let
$M_A$ and $M_B$ be matchings covering $A$ and $B$, with $|M_A|=|A|<|B|=|M_B|$. In $M_A\triangle M_B$, every vertex has degree at most $2$, so the components are paths and cycles alternating between the two
matchings, and since $M_B$ has more edges, some component is a path $P$ with one more edge of $M_B$ than of $M_A$, starting and ending with edges of $M_B$. Its two ends are on opposite sides; let $x$ be
the end in $E$. It is covered by $M_B$, so $x\in B$, and not by $M_A$ (an edge of $M_A$ at $x$ would belong to the component), so $x\notin A$. As in Berge's theorem (Theorem~\ref{thm:berge}), $M_A\triangle P$
is a matching, and it covers $A+x$.
\end{proof}

Graphic matroids are linear over every field (Exercise~\ref{ex:graphiclinear}), and the transversal ones are the matroids of matchings: their augmentation is Berge's augmenting path. Here is one
that comes from scheduling, the problem that best shows what a matroid can do for a programmer.

\begin{example}[scheduling with deadlines]\label{exa:scheduling}
A machine does one job per day. Job $j$ takes a day, must be done by day $d_j$, and earns $w_j>0$ if it is. A set of jobs is \emph{feasible} if all of them can be done on time. Assigning
job $j$ to a day in $S_j=\{1,\dots,d_j\}$, the feasible sets are the partial transversals of $(S_j)$, with the roles of elements and sets exchanged: the jobs are the elements, and each day may
receive the jobs whose set contains it. More simply, by Hall's theorem (Theorem~\ref{thm:hall}) or a direct argument, a set $A$ is feasible if and only if, for each $t$, at most $t$ jobs of
$A$ have deadline at most $t$; and $A$ is then done on time by doing its jobs in order of deadline. Either way the feasible sets form a matroid (they are the partial transversals of the
family of jobs $\{j:d_j\ge t\}$ indexed by the days $t$, or check (I3) directly in Exercise~\ref{ex:schedule}). By Theorem~\ref{thm:greedy} below, the best plan is greedy: consider the jobs by decreasing
reward and keep each one that leaves the set feasible.
\end{example}

\section{Bases, circuits and rank}

\begin{proposition}\label{prop:bases}
Let $M=(E,\mathcal I)$ be a matroid and $X\subseteq E$. All maximal independent subsets of $X$ have the same size, the \emph{rank} $r(X)$ of $X$. In particular all maximal independent sets, the
\emph{bases} of $M$, have the same size $r(M)=r(E)$, and every independent subset of $X$ extends to one of size $r(X)$.
\end{proposition}
\begin{proof}
If $A,B\subseteq X$ are independent and maximal among independent subsets of $X$, and $|A|<|B|$, then (I3) gives $x\in B\setminus A\subseteq X$ with $A+x$ independent, contradicting the maximality of $A$. An
independent subset of $X$ grows, one element at a time, into a maximal one.
\end{proof}

For a vector configuration, $r(X)$ is the dimension of the span of $X$, and the proposition is the theorem that all bases of a vector space have the same size. For a graph with $n$ vertices and $c$
components, $r(E)=n-c$, and all spanning forests have $n-c$ edges. A \emph{circuit} is a minimal dependent set: a cycle of a graph, a minimal linearly dependent family of vectors. A \emph{loop} is an
element $e$ with $\{e\}$ dependent, and two non-loops are \emph{parallel} if they form a circuit.

\begin{proposition}[circuits]\label{prop:circuits}
\begin{enumerate}[label=(\alph*)]
\item (\emph{Elimination}) If $C_1\ne C_2$ are circuits and $e\in C_1\cap C_2$, there is a circuit $C\subseteq(C_1\cup C_2)-e$.
\item (\emph{Unique circuit}) If $I$ is independent and $I+x$ is dependent, then $I+x$ contains exactly one circuit, and it contains $x$.
\end{enumerate}
\end{proposition}
\begin{proof}
(a) Suppose $(C_1\cup C_2)-e$ is independent, of size $|C_1\cup C_2|-1$. Since circuits are minimal and $C_1\ne C_2$, $C_1\not\subseteq C_2$; pick $f\in C_1\setminus C_2$. The set $C_1-f$ is independent;
extend it to a maximal independent subset $I$ of $C_1\cup C_2$. Then $f\notin I$, otherwise $C_1\subseteq I$; and some $g\in C_2$ is not in $I$, otherwise $C_2\subseteq I$; and $g\ne f$ as $f\notin C_2$. So
$|I|\le|C_1\cup C_2|-2$, while $(C_1\cup C_2)-e$ is an independent subset of $C_1\cup C_2$ of larger size. This contradicts Proposition~\ref{prop:bases}.

(b) A dependent set contains a circuit (take a minimal dependent subset). Every circuit in $I+x$ contains $x$, since $I$ is independent. If there were two, the elimination of $x$ would give a circuit
inside $I$.
\end{proof}

The rank function carries the whole matroid: $A$ is independent exactly when $r(A)=|A|$. Its key property is the discrete analogue of the dimension formula $\dim(U+V)+\dim(U\cap V)=\dim U+\dim V$.

\begin{theorem}[submodularity]\label{thm:submodular}
The rank function satisfies $0\le r(A)\le|A|$, $r(A)\le r(B)$ for $A\subseteq B$, and
\[
r(A\cup B)+r(A\cap B)\le r(A)+r(B).
\]
\end{theorem}
\begin{proof}
The first two properties are clear. Let $I$ be a maximal independent subset of $A\cap B$, and extend it to a maximal independent subset $J$ of $A\cup B$ (Proposition~\ref{prop:bases}), so $|I|=r(A\cap B)$
and $|J|=r(A\cup B)$. The set $J\cap A\cap B$ is independent, contains $I$, and lies in $A\cap B$, so it equals $I$. The sets $J\cap A$ and $J\cap B$ are independent subsets of $A$ and $B$, so
\[
r(A)+r(B)\ge|J\cap A|+|J\cap B|=|J|+|J\cap A\cap B|=r(A\cup B)+r(A\cap B).
\]
\end{proof}

Conversely, any function with these three properties is the rank function of a matroid (Exercise~\ref{ex:rankaxioms}). Matroids can also be defined by their bases, their circuits, or their closure
operator; Whitney's paper already proved that all these axiom systems are equivalent. We need the version with bases.

\begin{theorem}[bases]\label{thm:bases}
Let $\mathcal B$ be a non-empty family of subsets of a finite set $E$. It is the family of bases of a matroid (namely of $\mathcal I=\{A:A\subseteq B\text{ for some }B\in\mathcal B\}$) if and only if it satisfies the
\emph{exchange axiom}: for $B_1,B_2\in\mathcal B$ and $x\in B_1\setminus B_2$, there is $y\in B_2\setminus B_1$ such that $B_1-x+y\in\mathcal B$.
\end{theorem}
\begin{proof}
\emph{Bases satisfy the exchange axiom.} $B_1-x$ is independent and smaller than $B_2$, so by (I3) some $y\in B_2\setminus(B_1-x)$ makes $B_1-x+y$ independent; $y\neq x$ since $x\notin B_2$, so $y\in B_2\setminus B_1$, and
$B_1-x+y$ is independent of size $r(M)$, hence a basis.

\emph{The exchange axiom gives a matroid.} (I1) and (I2) hold for $\mathcal I$. First, all members of $\mathcal B$ have the same size. Otherwise choose $B_1,B_2\in\mathcal B$ with $|B_1|<|B_2|$ and $|B_2\setminus B_1|$
as small as possible, and $x\in B_2\setminus B_1$. The exchange axiom gives $y\in B_1\setminus B_2$ with $B_2'=B_2-x+y\in\mathcal B$. Then $|B_2'|=|B_2|>|B_1|$ and $|B_2'\setminus B_1|=|B_2\setminus B_1|-1$, a contradiction.

For (I3), let $A,C\in\mathcal I$ with $|A|<|C|$. Among all pairs of members $B_A\supseteq A$ and $B_C\supseteq C$ of $\mathcal B$, choose one minimising $s=|B_C\setminus(C\cup B_A)|+|B_A\setminus(A\cup B_C)|$. If some
$x\in C\setminus A$ lies in $B_A$, then $A+x\subseteq B_A$ is independent and we are done; assume not.

Suppose the first term of $s$ is positive, with $x\in B_C\setminus(C\cup B_A)$. Exchange gives $y\in B_A\setminus B_C$ with $B_C-x+y\in\mathcal B$; it contains $C$ (as $x\notin C$), the first term drops by $1$,
and the second cannot increase, since $B_C$ only gained $y$: a contradiction. Suppose the second term is positive, with $y\in B_A\setminus(A\cup B_C)$. Exchange gives $z\in B_C\setminus B_A$ with
$B_A-y+z\in\mathcal B$, containing $A$. If $z\in C$, then $z\in C\setminus A$ and $A+z\subseteq B_A-y+z$: done. Otherwise $z\notin C$, the second term drops by $1$, and the first cannot increase,
since $B_A$ lost only $y\notin B_C$: a contradiction again.

So $B_C\subseteq C\cup B_A$ and $B_A\subseteq A\cup B_C$. Let $T=B_A\cap B_C$. Then $B_A=A\cup(T\setminus A)$ and $B_C=C\cup(T\setminus C)$. As no element of $C\setminus A$ is in $B_A$, $T\cap C\subseteq A$, so
$T\setminus A\subseteq T\setminus C$. Hence $|B_C|=|C|+|T\setminus C|\ge|C|+|T\setminus A|>|A|+|T\setminus A|=|B_A|$, contradicting the equality of sizes.
\end{proof}

\section{The greedy algorithm}

Let $E$ carry a weight $w:E\to\R_{\ge0}$, and $w(A)=\sum_{a\in A}w(a)$. Given a family $\mathcal I$ satisfying (I1) and (I2), the \emph{greedy algorithm} sorts $E$ by decreasing weight,
$w(e_1)\ge w(e_2)\ge\dots\ge w(e_n)$ (ties broken arbitrarily), starts from $G=\varnothing$, and for $i=1,\dots,n$ replaces $G$ by $G+e_i$ if this set is in $\mathcal I$.

\begin{theorem}[Rado, 1957; Gale, 1968; Edmonds, 1971]\label{thm:greedy}
Let $\mathcal I$ be a family of subsets of $E$ satisfying (I1) and (I2). The greedy algorithm returns an independent set of maximum weight for every weight function $w\ge0$ if and only if $\mathcal I$ satisfies
(I3), that is, if and only if $(E,\mathcal I)$ is a matroid.
\end{theorem}
\begin{proof}
\emph{Matroids.} Let $G=\{g_1,\dots,g_k\}$ be the greedy output, in the order chosen, so $w(g_1)\ge\dots\ge w(g_k)$. It is a basis: an element $x$ rejected when the current set was $G'\subseteq G$ has $G'+x$
dependent, so $G+x$ is dependent by (I2). Since weights are non-negative, some basis $O$ has maximum weight; list it as $o_1,\dots,o_k$ with $w(o_1)\ge\dots\ge w(o_k)$. We show $w(g_i)\ge w(o_i)$ for all $i$, which gives
$w(G)\ge w(O)$.

Suppose $w(g_i)<w(o_i)$ for some $i$. The sets $A=\{g_1,\dots,g_{i-1}\}$ and $B=\{o_1,\dots,o_i\}$ are independent with $|A|<|B|$, so some $o_j\in B\setminus A$ has $A+o_j$ independent. Now $w(o_j)\ge w(o_i)>w(g_i)$, so the
algorithm examined $o_j$ before $g_i$, when its current set $G'$ was contained in $A$. Then $G'+o_j\subseteq A+o_j$ is independent, and the algorithm took $o_j$: $o_j\in A$, a contradiction.

\emph{Only matroids.} Suppose (I3) fails: $A,B\in\mathcal I$, $|A|=a<|B|=b$, and $A+x\notin\mathcal I$ for every $x\in B\setminus A$. Choose $0<\varepsilon<\frac1{a+1}$ and put $w=1+\varepsilon$ on $A$, $w=1$ on $B\setminus A$, $w=0$ elsewhere.
The algorithm first examines the elements of $A$ and takes them all, since subsets of $A$ are independent. Then every $x\in B\setminus A$ is rejected, since the current set contains $A$, and $A+x$ is dependent,
hence so is every set containing it. The remaining elements weigh $0$. So greedy returns weight $a(1+\varepsilon)<a+1\le b\le w(B)$, which is not the maximum.
\end{proof}

\begin{corollary}[Kruskal, 1956]\label{cor:kruskal}
Let $G$ be a connected graph with edge costs $c$. Examining the edges by increasing cost and keeping each edge that closes no cycle with those already kept yields a spanning tree of minimum total cost.
\end{corollary}
\begin{proof}
Spanning trees are the bases of $M(G)$, all with $n-1$ edges. With $C\ge\max c$, the weight $w=C-c\ge0$ gives a basis $T$ the weight $(n-1)C-c(T)$, so minimum cost means maximum weight, and decreasing
$w$ means increasing $c$. Kruskal's rule is the greedy algorithm for $w$.
\end{proof}

\begin{figure}[ht]
\centering
\fig{ch10-kruskal}
\caption{Kruskal's algorithm on a graph with $7$ vertices. The kept edges (thick, numbered in the order they are kept) form a spanning tree of cost $39$. The dashed edge of cost $6$ was rejected
because it closes the cycle shown with the kept edges of costs $3$, $4$ and $5$, which are all cheaper: this is the exchange argument at work.}
\label{fig:kruskal}
\end{figure}

Prim's algorithm, which grows a single tree from a root, and Bor\r{u}vka's, which grows all components at once, are also correct, by a lemma about cuts (Exercise~\ref{ex:prim}). How fast they run is a
question of data structures, treated in Tome~C2. Greedy fails for matchings, which do not form a matroid (Exercise~\ref{ex:greedyfail}): the family of matchings satisfies (I1) and (I2), and
the theorem then guarantees that some weights defeat greed. For travelling salesmen it fails badly; Tome~C3 explains why no fast algorithm is expected there at all.

\section{Duality and minors}

\begin{theorem}[duality, Whitney, 1935]\label{thm:dual}
Let $M$ be a matroid on $E$ with family of bases $\mathcal B$. Then $\mathcal B^*=\{E\setminus B:B\in\mathcal B\}$ is the family of bases of a matroid $M^*$, the \emph{dual} of $M$, with rank function
\[
r^*(A)=|A|-r(E)+r(E\setminus A),
\]
and $(M^*)^*=M$.
\end{theorem}
\begin{proof}
By Theorem~\ref{thm:bases} we check the exchange axiom for $\mathcal B^*$. Let $B_1^*=E\setminus B_1$, $B_2^*=E\setminus B_2$ and $x\in B_1^*\setminus B_2^*$, that is $x\in B_2\setminus B_1$. We need $y\in B_2^*\setminus B_1^*=B_1\setminus B_2$ such that
$B_1^*-x+y$ is in $\mathcal B^*$, that is $B_1+x-y\in\mathcal B$. By Proposition~\ref{prop:circuits}(b), $B_1+x$ contains a unique circuit $C$, and $x\in C$. Since $B_2$ is independent, $C\not\subseteq B_2$; take $y\in C\setminus B_2$;
then $y\ne x$, so $y\in B_1\setminus B_2$. The set $B_1+x-y$ contains no circuit, since every circuit in $B_1+x$ is $C$, which contains $y$. It is independent of size $r(M)$, hence a basis.

For the rank, $r^*(A)=\max_{B}|A\setminus B|=|A|-\min_B|A\cap B|$ over bases $B$ of $M$. For every basis, $B\setminus A$ is an independent subset of $E\setminus A$, so $|A\cap B|=r(E)-|B\setminus A|\ge r(E)-r(E\setminus A)$, with
equality when $B$ is obtained by extending a maximal independent subset of $E\setminus A$ to a basis. Finally, complementing twice gives back $\mathcal B$.
\end{proof}

The dual of a linear matroid is linear: if $M$ is represented by the columns of an $r\times n$ matrix $(I_r\mid A)$, then $M^*$ is represented by the columns of $(-A^{\mathsf T}\mid I_{n-r})$ (Exercise~\ref{ex:lineardual}).
Coding theorists know this as the passage from a generator matrix to a parity-check matrix (Tome~C4). For graphs, duality turns cycles into cuts. A \emph{bond} of a connected graph is a minimal set of edges
whose removal disconnects it.

\begin{proposition}[cographic matroids]\label{prop:bonds}
Let $G$ be connected. The circuits of $M(G)^*$ are the bonds of $G$.
\end{proposition}
\begin{proof}
A set $D$ is independent in $M^*$ when it is contained in the complement of a basis, that is, when $E\setminus D$ contains a spanning tree, that is, when $G-D$ is connected. So $D$ is dependent in $M^*$ when $G-D$ is disconnected,
and the minimal such sets are the bonds.
\end{proof}

\begin{remark}
For a connected plane graph $G$, the \emph{dual} graph $G^*$ has a vertex in each face of $G$ and an edge $e^*$ crossing each edge $e$, joining the faces on the two sides of $e$. Whitney proved in 1933 that a set
of edges is a cycle of $G$ exactly when the corresponding edges form a bond of $G^*$; equivalently $M(G^*)=M(G)^*$, and $T$ is a spanning tree of $G$ exactly when the edges $e^*$, $e\notin T$, form a spanning tree of
$G^*$ (Figure~\ref{fig:dualtree}). Counting edges of the two trees gives $m=(n-1)+(f-1)$, which is Euler's formula of Chapter~7 again. Whitney also proved the converse: $G$ is planar if and only if $M(G)^*$ is graphic. So the
dual matroid is the algebraic shadow of the dual map, and it exists for every graph, while dual graphs exist only for planar ones.
\end{remark}

\begin{figure}[ht]
\centering
\fig{ch10-dual}
\caption{A plane graph (blue, $6$ vertices, $9$ edges) and its dual (orange, one vertex per face, $5$ faces including the outer one). A spanning tree of the graph (thick blue) and the dual edges of the
$4$ remaining edges (thick orange), which form a spanning tree of the dual: $9=5+4$.}
\label{fig:dualtree}
\end{figure}

The two basic operations on matroids generalise deleting and contracting an edge of a graph. For $e\in E$, the \emph{deletion} $M\setminus e$ is the matroid on $E-e$ whose independent sets are those of $M$ not containing $e$.
If $e$ is not a loop, the \emph{contraction} $M/e$ is the matroid on $E-e$ whose independent sets are the $I\subseteq E-e$ with $I+e$ independent in $M$; if $e$ is a loop, $M/e=M\setminus e$. (Both are matroids: the
first is clear, and for the second (I3) follows from (I3) in $M$ applied to $I+e$ and $J+e$.) For a graph, $M(G)\setminus e=M(G-e)$ and $M(G)/e=M(G/e)$. A \emph{minor} of $M$ is a matroid obtained by a sequence of deletions and
contractions; duality exchanges them: $(M\setminus e)^*=M^*/e$ (Exercise~\ref{ex:minordual}). Many classes of matroids are characterised by forbidden minors, in the spirit of Wagner's form of Kuratowski's
theorem (Theorem~\ref{thm:kuratowski}): Tutte proved in 1958 that a matroid is representable over $\mathbb F_2$ if and only if it has no minor $U_{2,4}$, and that it is graphic if and only if it has no minor
among $U_{2,4}$, the Fano matroid $F_7$, its dual, and the duals $M(K_5)^*$, $M(K_{3,3})^*$.

\section{Matroids that are not like vectors}

Which matroids are linear? Over which fields? The most famous example lives in the smallest projective plane. The \emph{Fano matroid} $F_7$ has as ground set the seven non-zero vectors of $\mathbb F_2^3$,
\[
1=100,\ 2=010,\ 3=001,\ 4=110,\ 5=101,\ 6=011,\ 7=111,
\]
with the linear independence of $\mathbb F_2$. Its rank is $3$, every pair is independent, and the dependent triples are the seven \emph{lines} $\{x,y,x+y\}$:
\[
124,\ 135,\ 236,\ 167,\ 257,\ 347,\ 456.
\]

\begin{figure}[ht]
\centering
\fig{ch10-fano}
\caption{The Fano plane: seven points and seven lines (six segments and the red circle, which is the line $456$), each line with three points, any two points on exactly one line. As a matroid, the independent sets are the sets of at most three
points not forming a line.}
\label{fig:fano}
\end{figure}

\begin{theorem}\label{thm:fano}
The Fano matroid is representable over a field $F$ if and only if $F$ has characteristic $2$.
\end{theorem}
\begin{proof}
Over a field of characteristic $2$, the seven vectors above, with entries $0$ and $1$ read in $F$, have the same dependent triples (a $3\times3$ determinant with entries $0,1$ computed in $F$ is the image of the integer determinant. Of the $35$ triples, the six lines other than $456$ have integer determinant $0$, the line $456$ has $\det(110,101,011)=-2$, which vanishes in characteristic $2$, and the remaining $28$ triples have determinant $\pm1$, which never vanishes). So $F_7$ is representable.

Conversely, let $v_1,\dots,v_7\in F^r$ represent $F_7$. They span a space of dimension $3$, which we identify with $F^3$; applying the invertible linear map sending $v_1,v_2,v_3$ (independent, as $123$ is not a line) to
the standard basis, we may assume $v_1=e_1$, $v_2=e_2$, $v_3=e_3$. As $124$ is a line, $v_4$ lies in the span of $e_1,e_2$: $v_4=(a,b,0)$, with $a,b\ne0$ because $\{1,4\}$ and $\{2,4\}$ are independent. Likewise
$v_5=(c,0,d)$ and $v_6=(0,e,f)$ with $c,d,e,f\ne0$. Write $v_7=(x,y,z)$. The lines $347$, $257$, $167$ give
\[
\det(e_3,v_4,v_7)=\pm(ay-bx)=0,\quad\det(e_2,v_5,v_7)=\pm(cz-dx)=0,\quad\det(e_1,v_6,v_7)=\pm(ez-fy)=0.
\]
If $x=0$ then $y=z=0$ and $v_7=0$, impossible. So $y=\frac{bx}a$, $z=\frac{dx}c$, and the third equation gives $\frac{edx}c=\frac{fbx}a$, that is $ade=bcf$. The line $456$ gives
\[
\det\begin{pmatrix}a&b&0\\c&0&d\\0&e&f\end{pmatrix}=-ade-bcf=0.
\]
So $2bcf=0$, with $bcf\ne0$: $2=0$ in $F$.
\end{proof}

Changing the line $456$ into an independent triple gives the \emph{non-Fano matroid}, which is representable exactly over fields of characteristic different from $2$ (Exercise~\ref{ex:nonfano}). Their direct sum is a
matroid representable over no field at all. In fact almost all matroids are not representable: Nelson proved in 2018 that the proportion of representable matroids on $n$ elements tends to $0$. Rota's
conjecture, that for each finite field there are finitely many excluded minors for representability, was announced proved by Geelen, Gerards and Whittle in 2014, at the end of a programme of fifteen years
modelled on the Robertson--Seymour graph minor theory; the full write-up is still in progress.

\begin{remark}[matroid intersection]
Given two matroids $M_1$, $M_2$ on the same ground set, how large can a set independent in both be? Edmonds proved in 1970 that
\[
\max\{|I|:I\in\mathcal I_1\cap\mathcal I_2\}=\min_{A\subseteq E}\bigl(r_1(A)+r_2(E\setminus A)\bigr),
\]
with a polynomial algorithm by augmenting paths in an ``exchange graph''; Tome~C2 proves it by linear programming. For the two partition matroids on the edges of a bipartite graph whose parts are the stars of the
vertices on each side, with capacities $1$, common independent sets are matchings, and the formula is K\H onig's theorem (Theorem~\ref{thm:konig}). With one partition matroid and one arbitrary matroid it becomes
Rado's theorem (1942): sets $S_1,\dots,S_m$ have a transversal independent in $M$ if and only if $r(\bigcup_{i\in J}S_i)\ge|J|$ for all $J$, which is Hall's theorem when $M$ is free. The notebook runs the intersection
algorithm. Intersecting \emph{three} matroids is already NP-hard (it contains the Hamiltonian path problem): two is the limit of the good theory.
\end{remark}

\begin{beauty}
Three rules, about empty sets, subsets, and a small set learning from a large one, are enough to capture what vectors, forests, matchings and schedules have in common, and they are exactly what makes greed
pay. One proof shows that all bases of a vector space and all spanning trees of a graph have the same size; one duality turns cycles into cuts, generators into parity checks, a map into its dual map; and the
smallest projective plane, seven points on seven lines, knows whether $1+1$ is zero.
\end{beauty}

\section{Exercises}

\begin{exercise}[$\star$]\label{ex:graphiclinear}
Orient each edge $uv$ of a graph $G$ and give it the vector $e_u-e_v\in F^V$, over any field $F$. Show that a set of edges is a forest if and only if its vectors are linearly independent, so that $M(G)$ is
representable over every field.
\end{exercise}

\begin{exercise}[$\star\star$]\label{ex:rankaxioms}
Let $r$ be a function on the subsets of a finite set $E$ with $0\le r(A)\le|A|$, $r$ monotone, and $r$ submodular. Show that $\mathcal I=\{A:r(A)=|A|\}$ is a matroid with rank function $r$.
\end{exercise}

\begin{exercise}[$\star$]\label{ex:prim}
\emph{The cut property.} Let $G$ be connected with distinct edge costs, $S$ a set of vertices, $\varnothing\ne S\ne V$, and $e$ the cheapest edge between $S$ and $V\setminus S$. Show that $e$ belongs to the minimum spanning tree.
Deduce that Prim's algorithm (grow a tree from a root by repeatedly adding the cheapest edge leaving it) and Bor\r{u}vka's algorithm are correct.
\end{exercise}

\begin{exercise}[$\star$]\label{ex:greedyfail}
Show that the matchings of the path with three edges $a,b,c$ (in this order) do not form a matroid, and give weights for which greedy returns a matching of non-maximum weight.
\end{exercise}

\begin{exercise}[$\star$]\label{ex:schedule}
Prove directly that the feasible sets of jobs of Example~\ref{exa:scheduling} satisfy (I3). Then solve the instance with deadlines $2,1,2,1,3,3$ and rewards $10,9,8,7,6,5$ for jobs $1,\dots,6$.
\end{exercise}

\begin{exercise}[$\star\star$]\label{ex:lineardual}
Let $M$ be represented by the columns of the $r\times n$ matrix $(I_r\mid A)$ over a field. Show that $M^*$ is represented by the columns of $(-A^{\mathsf T}\mid I_{n-r})$.
\end{exercise}

\begin{exercise}[$\star\star$]\label{ex:minordual}
Show that $(M\setminus e)^*=M^*/e$ and $(M/e)^*=M^*\setminus e$.
\end{exercise}

\begin{exercise}[$\star$]\label{ex:u24}
Show that $U_{2,4}$ is not representable over $\mathbb F_2$, and that it is representable over every field with at least $3$ elements. Is it graphic?
\end{exercise}

\begin{exercise}[$\star\star$]\label{ex:nonfano}
Show that the non-Fano matroid (the seven points with the six lines other than $456$) is representable over $F$ if and only if the characteristic of $F$ is not $2$.
\end{exercise}

\begin{exercise}[$\star$]\label{ex:orthogonal}
Show that a circuit $C$ of $M$ and a circuit $D$ of $M^*$ never have exactly one common element. What does it say about cycles and bonds of a graph?
\end{exercise}

\begin{exercise}[$\star$]\label{ex:treegraph}
\emph{The tree graph.} Join two spanning trees of a connected graph when they differ by exchanging one edge. Show that this graph is connected, and that the distance between $T_1$ and $T_2$ is $|T_1\setminus T_2|$.
\end{exercise}

\section{Hints}

\begin{hint}{ex:graphiclinear}
A cycle gives a signed sum of vectors equal to $0$; conversely, a dependence relation restricted to a forest is zero at the leaves.
\end{hint}
\begin{hint}{ex:rankaxioms}
For (I3), if $r(A+x)=r(A)$ for all $x\in B\setminus A$, add the elements one by one with submodularity to get $r(A\cup B)=r(A)$.
\end{hint}
\begin{hint}{ex:prim}
If a spanning tree $T$ avoids $e$, the path in $T$ between the ends of $e$ crosses the cut by another edge.
\end{hint}
\begin{hint}{ex:greedyfail}
Two matchings of different sizes, $\{b\}$ and $\{a,c\}$.
\end{hint}
\begin{hint}{ex:schedule}
Use the criterion: for each $t$, at most $t$ jobs with deadline at most $t$.
\end{hint}
\begin{hint}{ex:lineardual}
Bases of $M^*$ are complements of bases of $M$; compare a square submatrix of $(I\mid A)$ with the complementary one of $(-A^{\mathsf T}\mid I)$.
\end{hint}
\begin{hint}{ex:minordual}
Compare bases on both sides; treat $e$ a loop or a coloop separately.
\end{hint}
\begin{hint}{ex:u24}
$\mathbb F_2^2$ has only three non-zero vectors.
\end{hint}
\begin{hint}{ex:nonfano}
Rerun the proof of Theorem~\ref{thm:fano}: now $ade+bcf\ne0$ is required.
\end{hint}
\begin{hint}{ex:orthogonal}
If $C\cap D=\{e\}$, extend $C-e$ to a basis avoiding $D-e$\dots
\end{hint}
\begin{hint}{ex:treegraph}
Exchange an edge of $T_1\setminus T_2$ for one of $T_2\setminus T_1$.
\end{hint}

\section{Solutions}

\begin{solution}{ex:graphiclinear}
\idea A cycle is a closed walk, and going around it the vectors telescope; a forest has leaves, which pin coefficients to zero.
\writeup If the edges contain a cycle $u_1u_2\cdots u_ku_1$, then $\sum_i\pm(e_{u_i}-e_{u_{i+1}})=0$, with the sign $+$ when the edge is oriented from $u_i$ to $u_{i+1}$: a non-trivial relation. Conversely suppose
$\sum_{e\in F}\lambda_ev_e=0$ for a forest $F$ with some $\lambda_e\ne0$, and keep only the edges with $\lambda_e\ne0$: a non-empty forest has a leaf $u$, incident to exactly one remaining edge $e$, and the coordinate $u$ of the relation
is $\pm\lambda_e=0$, a contradiction.
\end{solution}

\begin{solution}{ex:rankaxioms}
\idea Submodularity says that adding elements to a set never helps more than adding them to a smaller set; if no single element of $B$ raises the rank of $A$, then all of $B$ together does not either.
\writeup (I1): $0\le r(\varnothing)\le0$. (I2): if $r(B)=|B|$ and $A\subseteq B$, submodularity with $A$ and $B\setminus A$ (disjoint) gives $|B|=r(B)\le r(A)+r(B\setminus A)\le r(A)+|B\setminus A|$, so $r(A)\ge|A|$. (I3): first a lemma: if
$r(X+x)=r(X)$ for all $x\in Y$, then $r(X\cup Y)=r(X)$. By induction on $|Y|$: with $Y=Y'+y$ and $r(X\cup Y')=r(X)$, submodularity for $X\cup Y'$ and $X+y$ gives $r(X\cup Y)+r(X)\le r(X\cup Y')+r(X+y)=2r(X)$, and monotonicity gives
equality. Now if $A,B\in\mathcal I$, $|A|<|B|$, and no $x\in B\setminus A$ has $r(A+x)=|A|+1$, then $r(A+x)=r(A)$ for all of them (rank grows by at most $1$ when adding an element: $r(A+x)\le r(A)+r(\{x\})\le r(A)+1$),
so $r(A\cup B)=r(A)=|A|<|B|=r(B)$, contradicting monotonicity. Finally the rank of this matroid is $r$: a maximal $I\subseteq X$ with $r(I)=|I|$ has $r(I+x)=r(I)$ for all $x\in X\setminus I$, so by the lemma $r(X)=r(I)=|I|$.
\end{solution}

\begin{solution}{ex:prim}
\idea A spanning tree that avoids the cheapest crossing edge can be improved by an exchange.
\writeup Let $T$ be the minimum spanning tree (unique when costs are distinct: Kruskal's algorithm makes no choices) and suppose $e=uv\notin T$, $u\in S$, $v\notin S$. The path in $T$ from $u$ to $v$ starts in $S$ and ends outside,
so it has an edge $f$ crossing the cut. Then $T-f+e$ is a spanning tree (removing $f$ splits $T$ into two parts, and $e$ joins them) and costs less, since $c(e)<c(f)$: a contradiction. Prim's algorithm always adds the
cheapest edge leaving its current tree $S$, so it adds only edges of $T$, and after $n-1$ steps it has built $T$. Bor\r{u}vka's adds, for each current component, its cheapest leaving edge: all belong to $T$ by the cut property, so
adding them all at once creates no cycle.
\end{solution}

\begin{solution}{ex:greedyfail}
\idea Look for two matchings of different sizes where the small one cannot learn from the large one.
\writeup $\{b\}$ and $\{a,c\}$ are matchings, but $\{b,a\}$ and $\{b,c\}$ are not: (I3) fails. With $w(b)=3$, $w(a)=w(c)=2$, greedy takes $b$ and stops with weight $3$, while $\{a,c\}$ weighs $4$. The theorem's proof built
exactly these weights. Maximum weight matchings can still be found in polynomial time (Edmonds' blossom algorithm, or as a matroid intersection in the bipartite case), but not greedily.
\end{solution}

\begin{solution}{ex:schedule}
\idea A set is feasible when no deadline is overloaded; adding the job of the bigger set whose deadline is least loaded keeps all prefixes within capacity.
\writeup Let $N_A(t)$ be the number of jobs of $A$ with deadline at most $t$; $A$ is feasible iff $N_A(t)\le t$ for all $t$ (necessity is clear; sufficiency: do the jobs in order of deadline, the $k$-th job, on day $k$, has
deadline $d\ge k$, since otherwise the $k$ jobs up to it would give $N_A(d)\ge k>d$). Let $A,B$ be feasible, $|A|<|B|$. Let $t_0$ be the largest
$t$ with $N_A(t)\ge N_B(t)$, or $0$ if there is none (note $N_A(\infty)=|A|<|B|$). Then $B$ has more jobs than $A$ with deadline greater than $t_0$, so some such job $x\in B\setminus A$; take one with the smallest deadline
$d_x>t_0$. Adding $x$ to $A$ changes $N(t)$ only for $t\ge d_x$, and there $N_A(t)<N_B(t)\le t$, so $N_{A+x}(t)\le t$: $A+x$ is feasible.
The instance: greedy by reward takes job $1$ (deadline $2$), job $2$ (deadline $1$), rejects job $3$ (three jobs with deadline at most $2$) and job $4$, takes job $5$ (deadline $3$), rejects job $6$. Total reward $10+9+6=25$, with the
schedule $2,1,5$ on days $1,2,3$.
\end{solution}

\begin{solution}{ex:lineardual}
\idea A basis of $M$ is a set of columns of $(I\mid A)$ whose determinant is non-zero; the matching minor of the other matrix has the same determinant up to sign.
\writeup Index the columns of both matrices by $E=R\sqcup S$, $|R|=r$, $|S|=n-r$. For $X\subseteq E$ with $|X|=r$, write $X=(R\setminus R_1)\cup S_1$ with $R_1\subseteq R$, $S_1\subseteq S$, $|R_1|=|S_1|$. Expanding along the identity
columns, the columns $X$ of $(I\mid A)$ have determinant $\pm\det A[R_1,S_1]$ (the rows $R_1$ and columns $S_1$ of $A$). The complement $E\setminus X=R_1\cup(S\setminus S_1)$, in $(-A^{\mathsf T}\mid I)$, has determinant
$\pm\det(-A^{\mathsf T})[S_1,R_1]=\pm\det A[R_1,S_1]$. So $X$ is a basis of $M$ exactly when $E\setminus X$ is a basis of the matroid of $(-A^{\mathsf T}\mid I)$, which therefore is $M^*$. In coding theory, $(I\mid A)$
generates a code and $(-A^{\mathsf T}\mid I)$ is its parity-check matrix.
\end{solution}

\begin{solution}{ex:minordual}
\idea Write all four matroids through their bases and complement.
\writeup Suppose $e$ is neither a loop nor a coloop (an element in every basis) of $M$. Bases of $M\setminus e$: bases of $M$ avoiding $e$ (they exist, as $e$ is not a coloop). Bases of $M/e$: the sets $B-e$ for bases $B\ni e$.
Then the bases of $(M\setminus e)^*$ are $(E-e)\setminus B=(E\setminus B)-e$ for bases $B\not\ni e$, that is, $B^*-e$ for bases $B^*$ of $M^*$ containing $e$: the bases of $M^*/e$ ($e$ is not a loop of $M^*$, since it is not a
coloop of $M$). The second identity follows by applying the first to $M^*$ and dualising. If $e$ is a coloop of $M$, it is a loop of $M^*$; then $M\setminus e=M/e$ and $M^*\setminus e=M^*/e$, and one checks the bases
directly in the same way.
\end{solution}

\begin{solution}{ex:u24}
\idea In $U_{2,4}$ any two elements are independent and any three dependent: four pairwise non-parallel vectors in a plane.
\writeup A representation lies in a $2$-dimensional space, with $4$ pairwise independent vectors, so on $4$ different lines through the origin. $\mathbb F_q^2$ has $q+1$ lines, so this needs $q\ge3$; over $\mathbb F_2$ there are only three.
Over a field with at least $3$ elements, $(1,0),(0,1),(1,1),(1,a)$ with $a\notin\{0,1\}$ work, and over infinite fields $(1,t)$ for four values of $t$. It is not graphic, since graphic matroids are representable over $\mathbb F_2$
(Exercise~\ref{ex:graphiclinear}).
\end{solution}

\begin{solution}{ex:nonfano}
\idea Same coordinates as for the Fano plane; only the last equation changes from ``$=0$'' to ``$\ne0$''.
\writeup Following the proof of Theorem~\ref{thm:fano}, a representation can be normalised to $v_1,v_2,v_3$ the standard basis, $v_4=(a,b,0)$, $v_5=(c,0,d)$, $v_6=(0,e,f)$, $v_7=(1,\frac ba,\frac dc)$, with $ade=bcf$, and now $456$ independent:
$ade+bcf\ne0$, that is $2bcf\neq0$, so the characteristic is not $2$. Conversely, if $2\ne0$, take $a=b=c=d=e=f=1$, $v_7=(1,1,1)$: one checks that the dependent triples are the six lines $124,135,236,167,257,347$, and that
$\det(v_4,v_5,v_6)=-2\ne0$; the other $28$ triples have determinant $\pm1$ by the computation in the proof of Theorem~\ref{thm:fano} (the notebook checks all $35$).
\end{solution}

\begin{solution}{ex:orthogonal}
\idea A circuit of $M^*$ is a minimal set meeting every basis of $M$; one element of overlap would let a basis slip through.
\writeup By Theorem~\ref{thm:dual}, $D$ is dependent in $M^*$ iff it meets every basis of $M$, and a circuit $D$ of $M^*$ is a minimal such set (its complement $E\setminus D$ is then a \emph{hyperplane}: a maximal set of rank
$r(E)-1$). Suppose $C\cap D=\{e\}$. By minimality, $D-e$ misses some basis $B$, so $B\subseteq(E\setminus D)+e$; hence $r((E\setminus D)+e)=r(E)$ while $r(E\setminus D)\le r(E)-1$ (as $E\setminus D$ contains no basis): $e$ is not in the
span of $E\setminus D$. But $C-e\subseteq E\setminus D$, and $e$ is in the span of $C-e$ since $C$ is a circuit: a contradiction. For a graph, a cycle and a bond never share exactly one edge: a cycle that crosses a cut must
cross back.
\end{solution}

\begin{solution}{ex:treegraph}
\idea Each exchange step brings $T_1$ one edge closer to $T_2$, using the exchange axiom; no step can do better than one edge.
\writeup If $T_1\ne T_2$, take $x\in T_1\setminus T_2$; by Theorem~\ref{thm:bases} there is $y\in T_2\setminus T_1$ with $T_1-x+y$ a spanning tree, adjacent to $T_1$ and with $|(T_1-x+y)\setminus T_2|=|T_1\setminus T_2|-1$. By induction the distance is at most
$|T_1\setminus T_2|$, and it is at least that, since one exchange changes $|T\setminus T_2|$ by at most $1$. The same holds for the bases of any matroid; for graphs, the tree graph is the state space of the random
walks that sample spanning trees (Chapter~9's notebook).
\end{solution}

\section*{Sources and further reading}

\begin{itemize}
\item H.~Whitney, \emph{On the abstract properties of linear dependence}, American Journal of Mathematics 57 (1935), 509--533: the founding paper, still readable; T.~Nakasawa, \emph{Zur Axiomatik der linearen
Abh\"angigkeit I--III}, Science Reports of the Tokyo Bunrika Daigaku (1935--1936), rediscovered by Nishimura and Kuroda, \emph{A Lost Mathematician, Takeo Nakasawa}, Birkh\"auser, 2009.
\item O.~Bor\r{u}vka, \emph{O jist\'em probl\'emu minim\'aln\'\i m} (On a certain minimal problem), Pr\'ace Moravsk\'e P\v{r}\'\i rodov\v{e}deck\'e Spole\v{c}nosti 3 (1926), 37--58; J.~B.~Kruskal, \emph{On the shortest spanning
subtree of a graph and the traveling salesman problem}, Proceedings of the AMS 7 (1956), 48--50; R.~L.~Graham and P.~Hell, \emph{On the history of the minimum spanning tree problem}, Annals of the History of
Computing 7 (1985), 43--57.
\item R.~Rado, \emph{Note on independence functions}, Proceedings of the London Mathematical Society 7 (1957), 300--320; D.~Gale, \emph{Optimal assignments in an ordered set}, Journal of Combinatorial Theory 4
(1968), 176--180; J.~Edmonds, \emph{Matroids and the greedy algorithm}, Mathematical Programming 1 (1971), 127--136, and \emph{Submodular functions, matroids, and certain polyhedra}, in Combinatorial Structures
and their Applications, Gordon and Breach, 1970, 69--87.
\item J.~Oxley, \emph{Matroid Theory}, 2nd ed., Oxford University Press, 2011: the reference, from which the proof of Theorem~\ref{thm:bases} is adapted; D.~J.~A.~Welsh, \emph{Matroid Theory}, Academic Press, 1976
(Dover reprint 2010); W.~T.~Tutte, \emph{A homotopy theorem for matroids I, II}, Transactions of the AMS 88 (1958), 144--174.
\item J.~Geelen, B.~Gerards and G.~Whittle, \emph{Solving Rota's conjecture}, Notices of the AMS 61 (2014), 736--743: an account written for a general audience; P.~Nelson, \emph{Almost all matroids are
non-representable}, Bulletin of the London Mathematical Society 50 (2018), 245--248, a two-page research paper whose counting argument is accessible after this chapter.
\end{itemize}
